\chapter{Modul Wali Kelas (Homeroom Teacher)}

Modul Wali Kelas adalah antarmuka khusus yang dirancang untuk memenuhi kebutuhan sehari-hari seorang Wali Kelas (\textit{Teacher}) tanpa membebankan mereka dengan kompleksitas panel administrasi penuh. Modul ini terdiri dari satu halaman utama: \textbf{Rekap Kehadiran Hari Ini}, sebuah dasbor cepat yang menampilkan kondisi kehadiran kelas asuhan secara \textit{real-time} pada hari yang sedang berjalan.

\section{Konsep Dasar \& Matriks Peran}
\begin{itemize}
    \item \textbf{Wali Kelas (\textit{Teacher})}: Satu-satunya peran yang memiliki akses ke menu \textbf{Menu Wali Kelas} di kelompok navigasi \textit{Wali Kelas}. Wali Kelas hanya dapat melihat data kelas yang secara resmi didelegasikan kepada mereka di pengaturan \textit{Master Data} kelas.
    \item \textbf{Admin / Super Admin}: Juga dapat mengakses halaman ini jika diperlukan, namun menu ini tidak tertera di panel navigasi mereka secara bawaan.
\end{itemize}

\section{Skenario Dunia Nyata \& Alur Kerja}

Ibu Sari, Wali Kelas 11 IPA 3, membuka EduConnect di \textit{smartphone} pada jam 08:30 pagi. Fase 1: Ia mengklik menu \textbf{Menu Wali Kelas} di panel navigasi kiri. Halaman \textbf{Rekap Kehadiran Hari Ini} terbuka dan secara otomatis memilih kelas asuhan beliau (11 IPA 3) karena itu satu-satunya kelas yang didelegasikan. Fase 2: Ibu Sari melihat daftar siswa beserta status kehadiran mereka hari ini. Ia memfilter tampilannya hanya pada siswa berstatus \textbf{Alpha} untuk mengetahui siapa yang belum hadir dan belum ada keterangan. Fase 3: Ibu Sari menggunakan informasi ini untuk menghubungi orang tua siswa yang bersangkutan.

\begin{figure}[H]
    \centering
    \includegraphics[width=0.9\textwidth]{placeholder_homeroom_today.png}
    \caption{Halaman Rekap Kehadiran Hari Ini untuk Wali Kelas, dengan filter kelas dan status serta komponen \textit{live} Livewire.}
    \label{figure:9.1}
\end{figure}

\section{Panduan Penggunaan (Step-by-Step)}

\subsection{Mengakses Rekap Kehadiran Hari Ini}
\begin{enumerate}
    \item Setelah \textit{login} dengan akun Wali Kelas (\textit{Teacher}), cari kelompok menu \textbf{Wali Kelas} di bagian bawah panel navigasi kiri.
    \item Klik menu \textbf{Menu Wali Kelas}. Sistem akan membuka halaman \textbf{Rekap Kehadiran Hari Ini} di URL \texttt{/wali-kelas/rekap-hari-ini}.
    \item Jika Anda belum ditetapkan sebagai wali kelas aktif untuk kelas manapun, halaman akan menampilkan notifikasi kuning yang berbunyi \textit{``Anda belum ditetapkan sebagai wali kelas aktif.''} Hubungi Admin untuk memperbarui data kelas.
\end{enumerate}

\subsection{Memantau dan Menyaring Status Kehadiran}
\begin{enumerate}
    \item Pada halaman \textbf{Rekap Kehadiran Hari Ini}, perhatikan baris filter di bagian atas.
    \item \textit{Dropdown} pertama menampilkan daftar kelas asuhan Anda. Jika Anda hanya menjadi Wali Kelas untuk satu kelas, kelas tersebut sudah terpilih secara otomatis. Mengubah pilihan \textit{dropdown} ini akan langsung memuat ulang data tanpa perlu menekan tombol Filter.
    \item \textit{Dropdown} kedua adalah filter \textbf{Status}. Pilih salah satu dari: \textbf{Semua Status}, \textbf{Hadir}, \textbf{Terlambat}, \textbf{Izin}, \textbf{Sakit}, atau \textbf{Alpha}. Mengubah pilihan ini juga akan memuat ulang data secara otomatis.
    \item Pojok kanan baris filter menampilkan informasi tanggal saat ini yang tidak dapat diubah, karena halaman ini selalu menampilkan data \textbf{hari yang sedang berjalan} saja.
    \item Tabel kehadiran di bawahnya ditenagai oleh komponen \textit{Livewire} (\texttt{homeroom-bulk-attendance}) yang memperbarui data secara \textit{real-time}.
\end{enumerate}

\section{Penanganan Masalah (Edge Cases)}
\begin{itemize}
    \item \textbf{Dropdown Kelas Terisi namun Tabel Kosong}: Pastikan kelas yang dipilih memiliki siswa aktif yang terdaftar. Kelas baru yang belum diisi siswa akan menghasilkan tabel kosong.
    \item \textbf{Data Tidak Terbaru / Tidak Sinkron}: Komponen \textit{Livewire} bergantung pada sesi aktif dan koneksi \textit{server}. Jika data tampak tidak sinkron, coba \textit{refresh} halaman secara manual menggunakan tombol \textit{refresh} di \textit{browser}.
    \item \textbf{Tidak Ada Menu Wali Kelas di Panel Navigasi}: Pastikan akun Anda telah ditetapkan dengan peran \textbf{Teacher} (\textit{role = teacher}) oleh Super Admin melalui menu \textbf{Guru \& Staf}. Akun dengan peran \textit{Admin} atau \textit{Staff} tidak menampilkan menu Wali Kelas meskipun mungkin bisa mengakses URL-nya secara langsung.
\end{itemize}
